\chapter{\textbf{Validation, Tests et Déploiement}}
\vspace{-0.7cm}

% ─────────────────────────────────────────────
\section*{Introduction}
% ─────────────────────────────────────────────
\vspace{-0.6cm}
\noindent\rule{\textwidth}{0.2pt}
\vspace{0.05cm}
\large

Ce chapitre présente la démarche de validation adoptée pour la plateforme AquaFlow,
les tests mis en œuvre à différents niveaux, les résultats obtenus, l'architecture de déploiement
conteneurisée, ainsi que les pipelines d'intégration continue garantissant la
qualité du code à chaque évolution du projet.

\normalsize
\vspace{-0.5cm}

% ─────────────────────────────────────────────
\section{Stratégie de tests}
% ─────────────────────────────────────────────

La validation d'AquaFlow repose sur quatre niveaux complémentaires : des
\textbf{tests unitaires backend} isolant chaque service métier NestJS avec des dépendances
simulées, des \textbf{tests unitaires frontend} validant les hooks, slices Redux et composants
en isolation, des \textbf{tests unitaires Python} couvrant la logique des services Big Data
sans infrastructure active, des \textbf{tests d'intégration} exécutés automatiquement
en pipeline CI/CD avec de vrais services, et des \textbf{tests fonctionnels
manuels} couvrant les scénarios critiques de bout en bout, incluant le pipeline Big Data complet.
Le tableau~\ref{tab:strategie_tests} synthétise cette stratégie.

\begin{table}[H]
	\centering
	\renewcommand{\arraystretch}{1.3}
	\caption{Synthèse de la stratégie de tests AquaFlow}
	\begin{tabular}{|p{4cm}|p{4.8cm}|p{3.4cm}|p{2.2cm}|}
		\hline
		\textbf{Niveau} & \textbf{Outils} & \textbf{Périmètre} & \textbf{Volume} \\
		\hline
		Tests unitaires backend
		& Jest 30 + ts-jest + @nestjs/testing
		& 10 modules NestJS
		& 141 cas \\
		\hline
		Tests unitaires frontend
		& Jest + React Testing Library
		& Hooks + Redux slices + composants
		& 91 cas \\
		\hline
		Tests unitaires Python
		& pytest 9 + mocks
		& Services Big Data (pipeline, détection, agrégation)
		& 70 cas \\
		\hline
		Tests d'intégration (CI)
		& GitHub Actions + Docker Services
		& PostgreSQL~15, Redis~7
		& Automatisé \\
		\hline
		Tests fonctionnels
		& Navigateur + Postman
		& Scénarios complets + pipeline Big Data
		& 19 scénarios \\
		\hline
	\end{tabular}
	\label{tab:strategie_tests}
\end{table}

Au total, 302 cas de test ont été définis et exécutés, répartis sur les trois niveaux
automatisés (289 s'exécutent en intégration continue ; les 13 tests d'agrégation PySpark
nécessitent un environnement Spark local). Cette stratégie garantit une couverture complète
des fonctionnalités critiques de la plateforme, de la logique métier jusqu'aux scénarios
d'utilisation réels.

% ─────────────────────────────────────────────
\section{Tests unitaires}
% ─────────────────────────────────────────────

\subsection{Backend NestJS}

Les tests unitaires backend couvrent dix modules métier de l'application NestJS,
dont trois modules Big Data ajoutés lors du Sprint~6 (producteur Kafka, consommateur Kafka
et service analytique). Toutes les dépendances externes sont simulées par des mocks Jest,
permettant une exécution sans base de données ni infrastructure active.
Le tableau~\ref{tab:tests_backend} présente la répartition par module.

\begin{table}[H]
	\centering
	\renewcommand{\arraystretch}{1.3}
	\caption{Tests unitaires backend — répartition par module métier}
	\begin{tabular}{|p{6.5cm}|p{3cm}|c|c|}
		\hline
		\textbf{Module testé} & \textbf{Outil} & \textbf{Suites} & \textbf{Cas} \\
		\hline
		Authentification              & Jest & 6  & 13 \\
		\hline
		Gestion des alertes           & Jest & 4  & 10 \\
		\hline
		Moteur de workflows           & Jest & 8  & 19 \\
		\hline
		Communication IoT             & Jest & 4  & 14 \\
		\hline
		Notifications                 & Jest & 5  & 8  \\
		\hline
		Gestion des capteurs          & Jest & 7  & 18 \\
		\hline
		Gestion des stations          & Jest & 6  & 16 \\
		\hline
		Producteur Kafka \textit{(Big Data)} & Jest & 4 & 11 \\
		\hline
		Consommateur Kafka \textit{(Big Data)} & Jest & 7 & 16 \\
		\hline
		Service analytique \textit{(Big Data)} & Jest & 5 & 16 \\
		\hline
		\multicolumn{2}{|r|}{\textbf{Total}} & \textbf{56} & \textbf{141} \\
		\hline
	\end{tabular}
	\label{tab:tests_backend}
\end{table}

Les 141 cas couvrent notamment la gestion des tokens JWT, la détection des cycles
dans les workflows, le traitement des données capteurs, la connexion et publication
Kafka, la création des alertes d'anomalie et les méthodes analytiques
(\texttt{getOverview}, \texttt{getSensorStats}, \texttt{getKpis}).
La figure~\ref{fig:backend-tests} illustre les résultats d'exécution obtenus.

\begin{figure}[H]
	\centering
	\includegraphics[width=0.9\textwidth]{figures/backend-tests.png}
	\caption{Résultats d'exécution des 141 tests unitaires backend (Jest)}
	\label{fig:backend-tests}
\end{figure}

La couverture sur les services métier ciblés atteint des taux élevés,
comme illustré dans le tableau~\ref{tab:couverture_services}.

\begin{table}[H]
	\centering
	\renewcommand{\arraystretch}{1.3}
	\caption{Couverture de code par module métier (backend)}
	\begin{tabular}{|p{5cm}|c|c|c|c|}
		\hline
		\textbf{Module testé} & \textbf{Statements} & \textbf{Branches} & \textbf{Fonctions} & \textbf{Lignes} \\
		\hline
		Communication IoT         & 91,5\,\% & 63,6\,\% & 66,7\,\% & 91,1\,\% \\
		\hline
		Gestion des stations      & 92,3\,\% & 86,4\,\% & 100\,\%  & 100\,\%  \\
		\hline
		Moteur de workflows       & 95,7\,\% & 65,2\,\% & 88,9\,\% & 94,9\,\% \\
		\hline
		Gestion des alertes       & $\sim$85\,\% & $\sim$70\,\% & $\sim$80\,\% & $\sim$85\,\% \\
		\hline
		Authentification          & $\sim$80\,\% & $\sim$65\,\% & $\sim$75\,\% & $\sim$80\,\% \\
		\hline
		Gestion des capteurs      & 77,3\,\% & 41,5\,\% & 72,7\,\% & 79,1\,\% \\
		\hline
		Producteur Kafka          & $\sim$90\,\% & $\sim$80\,\% & $\sim$85\,\% & $\sim$90\,\% \\
		\hline
		Consommateur Kafka        & $\sim$85\,\% & $\sim$75\,\% & $\sim$80\,\% & $\sim$85\,\% \\
		\hline
		Service analytique        & $\sim$75\,\% & $\sim$60\,\% & $\sim$70\,\% & $\sim$75\,\% \\
		\hline
		Notifications             & 65,7\,\% & 43,3\,\% & 77,8\,\% & 68,3\,\% \\
		\hline
	\end{tabular}
	\label{tab:couverture_services}
\end{table}

Le taux de couverture dépasse 75\,\% pour la quasi-totalité des modules, atteignant
100\,\% pour la gestion des stations et 95,7\,\% pour le moteur de workflows.
La figure~\ref{fig:coverage-report} présente le rapport complet généré par Jest.

\begin{figure}[H]
	\centering
	\includegraphics[width=0.9\textwidth]{figures/coverage-report.png}
	\caption{Rapport de couverture de code Jest — backend NestJS}
	\label{fig:coverage-report}
\end{figure}

% ─────────────────────────────────────────────
\subsection{Frontend React}

Les tests unitaires frontend couvrent cinq modules clés de l'interface : le hook de
communication temps réel, les deux slices Redux gérant les alertes et les notifications,
la barre de navigation administrative et la page de détail d'un capteur.
Le tableau~\ref{tab:tests_frontend} en présente la répartition.

\begin{table}[H]
	\centering
	\renewcommand{\arraystretch}{1.3}
	\caption{Tests unitaires frontend — répartition par module}
	\begin{tabular}{|p{6.5cm}|p{3cm}|c|c|}
		\hline
		\textbf{Module testé} & \textbf{Outil} & \textbf{Suites} & \textbf{Cas} \\
		\hline
		Communication WebSocket   & Jest + RTL & 8  & 19 \\
		\hline
		Gestion des alertes       & Jest + RTL & 6  & 20 \\
		\hline
		Notifications             & Jest + RTL & 8  & 17 \\
		\hline
		Barre de navigation admin & Jest + RTL & —  & 11 \\
		\hline
		Page détail capteur       & Jest + RTL & —  & 24 \\
		\hline
		\multicolumn{2}{|r|}{\textbf{Total}} & \textbf{22+} & \textbf{91} \\
		\hline
	\end{tabular}
	\label{tab:tests_frontend}
\end{table}

Les 91 cas valident notamment la gestion des connexions et déconnexions WebSocket,
les transitions d'état Redux, le rendu conditionnel des composants et le traitement
des événements temps réel. La figure~\ref{fig:frontend-tests} présente les résultats
d'exécution obtenus.

\begin{figure}[H]
	\centering
	\includegraphics[width=0.9\textwidth]{figures/frontend-tests.png}
	\caption{Résultats d'exécution des 91 tests unitaires frontend (Jest + React Testing Library)}
	\label{fig:frontend-tests}
\end{figure}

% ─────────────────────────────────────────────
\subsection{Services Python — Module Big Data}

Les services Big Data sont validés par 70 cas de test Python répartis sur trois modules
fonctionnels. Ces tests s'appuient sur le framework pytest avec des objets simulés
reproduisant le comportement des services externes (Kafka, MinIO, Spark), sans nécessiter
de cluster actif. Les 57 premiers cas s'exécutent intégralement en pipeline CI/CD ;
les 13 tests d'agrégation PySpark nécessitent un environnement Spark local
(\texttt{pip install pyspark}). Le tableau~\ref{tab:tests_python} détaille cette répartition.

\begin{table}[H]
	\centering
	\renewcommand{\arraystretch}{1.3}
	\caption{Tests unitaires Python — répartition par module Big Data}
	\begin{tabular}{|p{4cm}|p{5.5cm}|c|p{2.5cm}|}
		\hline
		\textbf{Module testé} & \textbf{Périmètre couvert} & \textbf{Cas} & \textbf{Exécution} \\
		\hline
		Pipeline d'archivage
		& Logique de batching, construction des clés de partition
		Hive, conversion en fichiers Parquet
		& 28 & CI + local \\
		\hline
		Détection d'anomalies
		& Formule du z-score, seuil de détection, mappage des niveaux
		de sévérité (WARNING / ERROR / CRITICAL), schéma Spark
		& 29 & CI + local \\
		\hline
		Agrégation des KPIs
		& Agrégation horaire et journalière, calcul du score de santé
		par station \textit{(requiert PySpark)}
		& 13 & Local uniquement \\
		\hline
		\multicolumn{2}{|r|}{\textbf{Total}} & \textbf{70} & 57 en CI \\
		\hline
	\end{tabular}
	\label{tab:tests_python}
\end{table}

Ces tests garantissent la fiabilité du traitement des données capteurs, de la détection
des anomalies et de l'agrégation des métriques avant leur transmission au backend.
La figure~\ref{fig:python-tests} présente les résultats d'exécution pytest.

\begin{figure}[H]
	\centering
	\includegraphics[width=0.9\textwidth]{figures/python-tests.png}
	\caption{Résultats d'exécution des tests Python (pytest) — 57 tests CI, 0 échec}
	\label{fig:python-tests}
\end{figure}

% ─────────────────────────────────────────────
\section{Tests d'intégration et fonctionnels}
% ─────────────────────────────────────────────

\subsection{Tests d'intégration}

Les tests d'intégration sont assurés par les pipelines CI/CD GitHub Actions, qui
provisionnent automatiquement des instances TimescaleDB~2.14 et Redis~7 lors de chaque
exécution. Trois périmètres sont couverts : les endpoints REST critiques validés
via Supertest, la communication WebSocket vérifiée via les outils du navigateur,
et le flux IoT testé de bout en bout par publication MQTT sur le topic \texttt{sensors}
jusqu'à la diffusion Socket.IO.

\subsection{Tests fonctionnels manuels}

Dix-neuf scénarios de validation manuelle ont été conduits sur l'interface web et
via Postman, couvrant six domaines : l'authentification, le moteur de workflows,
la gestion des alertes, la supervision temps réel, la maintenance et le pipeline
Big Data complet. Le tableau~\ref{tab:tests_fonctionnels} présente les huit scénarios
applicatifs représentatifs.

\begin{table}[H]
	\centering
	\renewcommand{\arraystretch}{1.3}
	\caption{Scénarios de tests fonctionnels — modules applicatifs}
	\begin{tabular}{|c|p{2.5cm}|p{5.8cm}|p{3cm}|c|}
		\hline
		\textbf{ID} & \textbf{Module} & \textbf{Scénario} & \textbf{Résultat attendu} & \textbf{Statut} \\
		\hline
		T01 & Auth             & Connexion avec identifiants valides               & JWT retourné, dashboard affiché    & \textbf{\checkmark} \\
		\hline
		T02 & Auth             & Accès page protégée sans jeton                    & Redirection vers login             & \textbf{\checkmark} \\
		\hline
		T03 & Workflow Builder & Créer et exécuter un workflow linéaire            & Statut \texttt{COMPLETED} retourné & \textbf{\checkmark} \\
		\hline
		T04 & Workflow Builder & Workflow contenant un cycle                       & Erreur : cycle détecté             & \textbf{\checkmark} \\
		\hline
		T05 & IoT              & Capteur dépasse le seuil maximal                  & Alerte + notification temps réel   & \textbf{\checkmark} \\
		\hline
		T06 & Alertes          & Acquitter puis résoudre une alerte                & Transitions \texttt{acknowledged} $\to$ \texttt{resolved} & \textbf{\checkmark} \\
		\hline
		T07 & Maintenance      & Créer un ordre et assigner un technicien          & Notification envoyée au technicien & \textbf{\checkmark} \\
		\hline
		T08 & Dashboard        & Chargement des indicateurs KPI temps réel         & Données agrégées affichées         & \textbf{\checkmark} \\
		\hline
	\end{tabular}
	\label{tab:tests_fonctionnels}
\end{table}

Les cinq scénarios dédiés au pipeline Big Data sont présentés dans le
tableau~\ref{tab:tests_bigdata}.

\begin{table}[H]
	\centering
	\renewcommand{\arraystretch}{1.4}
	\caption{Scénarios de tests fonctionnels — pipeline Big Data (bout en bout)}
	\begin{tabular}{|c|p{3.5cm}|p{5.5cm}|p{2.5cm}|c|}
		\hline
		\textbf{ID} & \textbf{Scénario} & \textbf{Description} & \textbf{Résultat attendu} & \textbf{Statut} \\
		\hline
		T09
		& Ingestion capteur vers archive
		& Une mesure MQTT est reçue, publiée dans Kafka, consommée
		par le service d'archivage et écrite en fichier Parquet
		dans le data lake MinIO
		& Fichier Parquet présent dans MinIO avec la partition
		Hive correcte
		& \textbf{\checkmark} \\
		\hline
		T10
		& Agrégation KPI vers base de données
		& Le job d'agrégation lit les données brutes depuis MinIO,
		calcule les indicateurs par capteur et les écrit
		dans TimescaleDB
		& Lignes présentes dans \texttt{sensor\_aggregates} ;
		API analytique retourne des données
		& \textbf{\checkmark} \\
		\hline
		T11
		& Détection d'anomalie en temps réel
		& Spark Streaming détecte un z-score $\geq$ 2,5 sur une
		fenêtre glissante et publie un événement d'anomalie
		dans Kafka
		& Message d'anomalie présent sur le topic
		\texttt{sensors.anomalies} avec les champs
		z-score, moyenne et écart-type
		& \textbf{\checkmark} \\
		\hline
		T12
		& Alerte anomalie vers tableau de bord
		& Le consommateur NestJS reçoit l'événement d'anomalie,
		crée une alerte de type \texttt{ANOMALY} et la diffuse
		via Socket.IO
		& Alerte persistée en base ; frontend reçoit
		l'événement en temps réel
		& \textbf{\checkmark} \\
		\hline
		T13
		& Consultation des KPIs pré-calculés
		& L'endpoint \texttt{GET /api/analytics/kpis} retourne
		les indicateurs agrégés depuis \texttt{sensor\_aggregates}
		sans recalcul
		& Réponse HTTP~200 avec données structurées ;
		délai de réponse inférieur à 100~ms
		& \textbf{\checkmark} \\
		\hline
	\end{tabular}
	\label{tab:tests_bigdata}
\end{table}

L'ensemble des 19 scénarios fonctionnels s'est conclu avec succès, confirmant la
conformité de la plateforme aux exigences métier définies en phase de conception.

% ─────────────────────────────────────────────
\section{Analyse des résultats}
% ─────────────────────────────────────────────

L'ensemble de la campagne de tests aboutit à un taux de réussite de \textbf{100\,\%}
sur les 302 cas automatisés. Le tableau~\ref{tab:bilan_tests} présente le bilan global
par catégorie.

\begin{table}[H]
	\centering
	\renewcommand{\arraystretch}{1.3}
	\caption{Bilan global des tests AquaFlow}
	\begin{tabular}{|p{4.8cm}|c|c|p{4.5cm}|}
		\hline
		\textbf{Catégorie} & \textbf{Fichiers} & \textbf{Cas} & \textbf{Périmètre} \\
		\hline
		Tests unitaires backend
		& 10 & 141
		& 10 modules NestJS (dont 3 Big Data) \\
		\hline
		Tests unitaires frontend
		& 5 & 91
		& Hooks + Redux + Composants \\
		\hline
		Tests unitaires Python Big Data
		& 3 & 70
		& Pipeline, détection, agrégation \\
		\hline
		Tests fonctionnels manuels
		& — & 19 scénarios
		& Modules critiques + pipeline Big Data \\
		\hline
		\textbf{Total automatisé}
		& \textbf{18} & \textbf{302}
		& (289 en CI, 13 PySpark local) \\
		\hline
	\end{tabular}
	\label{tab:bilan_tests}
\end{table}

Les dix-huit modules couverts regroupent l'ensemble des services critiques de la
plateforme. Les services Big Data sont validés en isolation par les tests Python
et par cinq scénarios fonctionnels de bout en bout. Le taux de couverture dépasse
75\,\% pour la quasi-totalité des modules, atteignant 100\,\% pour la gestion des
stations et 95,7\,\% pour le moteur de workflows.

% ─────────────────────────────────────────────
\section{Déploiement de la plateforme}
% ─────────────────────────────────────────────

\subsection{Services et configuration}

La plateforme AquaFlow est conteneurisée dans un fichier Docker Compose unique
regroupant douze services organisés en trois couches : cinq services applicatifs,
quatre services Big Data et trois services d'infrastructure de support.
Le tableau~\ref{tab:services_docker} décrit chaque service, son image et son rôle.

\begin{table}[H]
	\centering
	\renewcommand{\arraystretch}{1.3}
	\caption{Services Docker Compose — images, ports et rôles}
	\begin{tabular}{|p{2.5cm}|p{3.8cm}|p{1.8cm}|p{4.8cm}|}
		\hline
		\textbf{Service} & \textbf{Image} & \textbf{Port} & \textbf{Rôle} \\
		\hline
		\texttt{frontend}
		& nginx:1.25-alpine
		& 3000 / 80
		& SPA React — interface opérateur \\
		\hline
		\texttt{backend}
		& node:20-alpine
		& 3001
		& API REST + WebSocket + intégration Kafka \\
		\hline
		\texttt{postgres}
		& timescale/timescaledb: 2.14.2-pg15
		& 5432
		& Séries temporelles + hypertables + agrégats continus \\
		\hline
		\texttt{redis}
		& redis:7-alpine
		& 6379
		& Cache sessions + liste de révocation JWT \\
		\hline
		\texttt{mosquitto}
		& eclipse-mosquitto:2
		& 1883, 9001
		& Broker MQTT — réception des mesures terrain \\
		\hline
		\texttt{kafka}
		& apache/kafka:3.7.1 (KRaft)
		& 9092
		& Bus de messages — ingestion des flux capteurs \\
		\hline
		\texttt{minio}
		& minio/minio:latest
		& 9000, 9002
		& Data lake S3-compatible — archivage Parquet \\
		\hline
		\texttt{minio-init}
		& minio/mc:latest
		& —
		& Initialisation unique des buckets MinIO \\
		\hline
		\texttt{kafka-to-minio}
		& build ./data-pipeline
		& —
		& Archivage continu Kafka $\to$ MinIO \\
		\hline
		\texttt{spark-master}
		& apache/spark:3.5.3
		& 8080, 7077
		& Nœud maître du cluster Spark \\
		\hline
		\texttt{spark-worker}
		& apache/spark:3.5.3
		& —
		& Nœud de calcul (1~Go RAM, 2~cœurs) \\
		\hline
		\texttt{spark-anomaly}
		& apache/spark:3.5.3
		& —
		& Détection d'anomalies Spark Structured Streaming \\
		\hline
	\end{tabular}
	\label{tab:services_docker}
\end{table}

Six services disposent d'un health check actif garantissant leur disponibilité avant
le démarrage des services dépendants : \texttt{postgres} (\texttt{pg\_isready}),
\texttt{redis} (\texttt{redis-cli ping}), \texttt{mosquitto} (publication test),
\texttt{kafka} (listage des topics), \texttt{minio} (endpoint de santé HTTP)
et \texttt{backend} (endpoint \texttt{/api/health}).

En production, seuls les ports 80, 443 et 9001 sont exposés vers l'extérieur,
limitant ainsi la surface d'attaque de l'infrastructure. Les services Big Data
opèrent entièrement en réseau interne Docker. Les builds utilisent des
\textit{multi-stage builds} séparant l'environnement de compilation de l'image
d'exécution, réduisant la taille des images finales et excluant les outils de
développement. Le backend s'exécute sous un utilisateur non-root dédié
(\texttt{aquaflow}).

La configuration repose sur des variables d'environnement injectées via un fichier
\texttt{.env} (modèle documenté : \texttt{.env.example}). Les variables obligatoires
couvrent les secrets JWT (\texttt{JWT\_SECRET}, \texttt{JWT\_REFRESH\_SECRET}),
les mots de passe des services, les URLs de l'application et les paramètres Big Data
(\texttt{KAFKA\_BROKERS}, \texttt{MINIO\_ENDPOINT},
\texttt{MINIO\_ACCESS\_KEY}, \texttt{MINIO\_SECRET\_KEY}).

\subsection{Procédure de déploiement}

La mise en œuvre nécessite Docker $\geq$~24.0 et Docker Compose $\geq$~2.20.
La procédure se résume à quatre étapes : clonage du dépôt, configuration du
fichier \texttt{.env}, lancement de la pile complète via
\texttt{docker-compose up --build -d}, puis initialisation de la base de données
lors du premier déploiement
(\texttt{docker exec aquaflow-backend npm run seed}).
Docker Compose gère l'ordre de démarrage via \texttt{depends\_on} conditionné au
statut \texttt{healthy} de chaque service — le backend attend notamment la disponibilité
de cinq services (\texttt{postgres}, \texttt{redis}, \texttt{mosquitto}, \texttt{kafka}
et \texttt{minio}) — garantissant un démarrage ordonné de l'ensemble de l'infrastructure
en une seule commande.

% ─────────────────────────────────────────────
\section{Pipeline CI/CD}
% ─────────────────────────────────────────────

La qualité du code est maintenue par trois pipelines GitHub Actions déclenchés
automatiquement sur chaque \textit{push} ou \textit{pull request} vers les branches
\texttt{main}, \texttt{master} et \texttt{develop}.
Chaque pipeline est déclenché sélectivement selon le répertoire modifié
(path filtering), évitant ainsi les exécutions inutiles lors de changements isolés.
Le tableau~\ref{tab:cicd_pipelines} présente la structure de chaque pipeline.

\begin{table}[H]
	\centering
	\renewcommand{\arraystretch}{1.3}
	\caption{Pipelines CI/CD GitHub Actions}
	\begin{tabular}{|p{2.8cm}|p{4cm}|p{4.5cm}|p{2.5cm}|}
		\hline
		\textbf{Pipeline} & \textbf{Jobs} & \textbf{Services provisionnés} & \textbf{Artefact} \\
		\hline
		\textbf{Backend CI}
		& 1.~Lint ESLint + build TypeScript \newline 2.~Tests unitaires + couverture
		& TimescaleDB~2.14-pg15 \newline Redis~7-alpine
		& Rapport couverture (14~j) \\
		\hline
		\textbf{Frontend CI}
		& 1.~Build production (\texttt{CI=true}) \newline 2.~Tests Jest
		& Aucun
		& Bundle React (7~j) \\
		\hline
		\textbf{Python CI}
		& 1.~Lint Flake8 \newline 2.~Tests Pytest + couverture
		& Aucun
		& Rapport couverture (14~j) \\
		\hline
	\end{tabular}
	\label{tab:cicd_pipelines}
\end{table}

Le pipeline \textbf{Backend CI} enchaîne un job de lint (ESLint + compilation TypeScript)
puis un job de tests avec provisionnement automatique de TimescaleDB~2.14 et Redis~7,
le rapport de couverture Jest étant archivé pendant 14~jours. Le pipeline
\textbf{Frontend CI} réalise un build de production avec \texttt{CI=true} traitant
les avertissements comme des erreurs, suivi des tests Jest. Le pipeline
\textbf{Python CI} exécute un lint Flake8 sur les services Big Data puis lance les
57 tests pytest sur les modules de traitement et de détection d'anomalies (les tests
d'agrégation PySpark sont exclus du pipeline CI et exécutés localement).
La figure~\ref{fig:ci-cd-pipeline} illustre l'exécution réussie des trois pipelines.

\begin{figure}[H]
	\centering
	\includegraphics[width=1\textwidth]{figures/ci-cd-pipeline.png}
	\caption{Pipelines CI/CD GitHub Actions — trois workflows en succès}
	\label{fig:ci-cd-pipeline}
\end{figure}

% ─────────────────────────────────────────────
\section*{Conclusion}
% ─────────────────────────────────────────────
\vspace{-0.6cm}
\noindent\rule{\textwidth}{0.2pt}
\vspace{0.05cm}
\large

Ce chapitre a présenté la démarche de validation et le déploiement de la plateforme
AquaFlow. Les 302 cas de test automatisés — répartis entre le backend NestJS (141),
le frontend React (91) et les services Python Big Data (70, dont 57 exécutés en
intégration continue) —, complétés par 19 scénarios fonctionnels dont 5 couvrant
le pipeline Big Data de bout en bout, attestent d'un taux de réussite de 100\,\%
sur l'ensemble des modules critiques du système.
L'infrastructure Docker Compose à douze services garantit un déploiement
reproductible et sécurisé en une seule commande, tandis que les trois pipelines
CI/CD assurent la qualité continue du code à chaque évolution du projet.

% ============================================================
